% ch4_b 第4章 §4.2 (书页 155–165 / p170–p180)

由于 $\langle T(c_{\pmb{k}}^{\dagger}(\tau)c_{\pmb{k}}(\tau'))\rangle = -\langle T(c_{\pmb{k}}(\tau')c_{\pmb{k}}^{\dagger}(\tau))\rangle$，我们有 $\langle T(c_{\omega_{\gamma}\pmb{k}}^{\dagger}c_{\omega_{\gamma}\pmb{k}})\rangle = -\langle T(c_{\omega_{\gamma}\pmb{k}}c_{\omega_{\gamma}\pmb{k}}^{\dagger})\rangle$。类似地，可以证明，对于实时间，
\begin{equation*}
\left\langle T(c_{\omega\pmb{k}}c_{\omega'\pmb{k}}^{\dagger})\right\rangle
= -\left\langle T(c_{\omega\pmb{k}}^{\dagger}c_{\omega'\pmb{k}})\right\rangle
= \mathrm{i}G^{\beta}(\omega,\pmb{k})(2\pi)^{-1}\delta(\omega-\omega')
\end{equation*}

这使我们能够直接在 $\omega$--$\pmb{k}$ 空间中进行编时关联的计算。我们还注意到，在编时平均（time-ordered average）中，$c$ 与 $c^{\dagger}$ 的表现如同反对易数（anti-commuting numbers）。

上面我们讨论了两个算符的编时关联。那么，多个算符的编时关联如何计算？这里我们有费米子的 Wick 定理。设 $O_{i}$ 是 $c$ 与 $c^{\dagger}$ 的线性组合。对于 $W = O_{1}O_{2}\dots O_{n}$，我们有
\begin{align*}
W = {}& {:W:} + (-)^{j_{1}-i_{1}-1}\sum_{(i_{1},j_{1})}{:W_{i_{1},j_{1}}:}\,\langle 0|O_{i_{1}}O_{j_{1}}|0\rangle\\
& \pm \sum_{\substack{(i_{1},j_{1}),(i_{2},j_{2})\\(i_{1},j_{1})\neq(i_{2},j_{2})}}{:W_{i_{1},j_{1},i_{2},j_{2}}:}\,\langle 0|O_{i_{1}}O_{j_{1}}|0\rangle\langle 0|O_{i_{2}}O_{j_{2}}|0\rangle + \dots \tag{4.2.11}
\end{align*}
其中 $(i_{1}, j_{1})$ 是满足 $i_{1} < j_{1}$ 的有序对，而 $\pm$ 由所需置换数目的奇偶性决定——这些置换把 $O_{j_{1}}$ 移到 $O_{i_{1}}$ 的右侧、把 $O_{j_{2}}$ 移到 $O_{i_{2}}$ 的右侧。由于按定义 $\langle 0|{:W:}|0\rangle = 0$，我们得到
\begin{equation*}
\langle 0|O_{1}O_{2}\dots O_{n}|0\rangle = \sum(-)^{P}\langle 0|O_{i_{1}}O_{i_{2}}|0\rangle\dots\langle 0|O_{i_{n-1}}O_{i_{n}}|0\rangle
\end{equation*}
其中 $\sum$ 遍历把 $1, 2, \dots, n$ 分组为有序对 $(i_{1}, i_{2})$、$(i_{3}, i_{4})$、$\dots$ 的所有可能方式（即 $i_{1} < i_{2}$，$i_{3} < i_{4}$，等等）。这里，若 $1, 2, \dots, n$ 与 $i_{1}, i_{2}, \dots, i_{n}$ 相差偶数个置换，则 $(-)^{P} = 1$；若 $1, 2, \dots, n$ 与 $i_{1}, i_{2}, \dots, i_{n}$ 相差奇数个置换，则 $(-)^{P} = -1$。例如，
\begin{gather*}
\langle 0|O_{1}O_{2}O_{3}O_{4}|0\rangle\\
= \langle 0|O_{1}O_{2}|0\rangle\langle 0|O_{3}O_{4}|0\rangle - \langle 0|O_{1}O_{3}|0\rangle\langle 0|O_{2}O_{4}|0\rangle + \langle 0|O_{1}O_{4}|0\rangle\langle 0|O_{2}O_{3}|0\rangle
\end{gather*}
应用于编时关联时，我们有
\begin{equation*}
\langle 0|T(O_{1}O_{2}\dots O_{n})|0\rangle = \sum(-)^{P}\langle 0|T(O_{i_{1}}O_{i_{2}})|0\rangle\dots\langle 0|T(O_{i_{n-1}}O_{i_{n}})|0\rangle
\end{equation*}
Wick 定理使我们能够用两算符关联来表示多算符关联。Wick 定理还蕴含 $\langle 0|T(\dots O_{m}\dots O_{n}\dots)|0\rangle \allowbreak= -\langle 0|T(\dots O_{n}\dots O_{m}\dots)|0\rangle$。因此，在多个 $c$ 与 $c^{\dagger}$ 算符的编时平均中（无论实时还是虚时），$c$ 与 $c^{\dagger}$ 的表现都如同反对易数。这一事实使我们能够为费米子系统构造路径积分表述（见 5.4.1 节）。

%FIG{4.4}: {“费米点”附近的态密度按 $N(\xi) \propto |\xi|^{d-1}$ 的方式趋于零。}

\subsection*{问题 4.2.1}

证明式 (4.2.7) 与式 (4.2.9)。

\subsection*{问题 4.2.2}

由式 (4.2.6) 证明式 (4.2.8)。

\subsection*{问题 4.2.3}

证明式 (4.2.10)。

\subsection*{问题 4.2.4}

就 $W = c_{k_{1}}c_{k_{2}}^{\dagger}c_{k_{3}}c_{k_{4}}^{\dagger}$ 与 $W = c_{k_{1}}c_{k_{2}}c_{k_{3}}^{\dagger}c_{k_{4}}^{\dagger}$ 两种情形，验证 Wick 定理的正确性。

\subsection{等空间格林函数与隧穿}

\begin{keypoints}
\item 格林函数在长时间上的代数衰减与跨费米面的无能隙激发相关。
\end{keypoints}

实时空中的格林函数由 $\mathrm{i}G(t, \pmb{x}_{1}, \pmb{x}_{2}) = \langle T(c(t, \pmb{x}_{1})c(0, \pmb{x}_{2})\rangle$ 给出。对于平移不变的系统（例如自由费米子系统），格林函数只依赖于 $\pmb{x}_{1} - \pmb{x}_{2}$，即 $G(t, \pmb{x}_{1}, \pmb{x}_{2}) = G(t, \pmb{x}_{1} - \pmb{x}_{2})$。对于自由费米子系统，我们有 $G(t, \pmb{x}) = \int \frac{\mathrm{d}^{d}\pmb{k}}{(2\pi)^{d}}G(t, \pmb{k})$。当 $\pmb{x} = 0$ 时（或者说，当 $\langle T(c(\pmb{x}_{1}, t_{1})c^{\dagger}(\pmb{x}_{2}, t_{2}))\rangle$ 中的费米算符位于等空间点（equal-space point）$\pmb{x}_{1} = \pmb{x}_{2}$ 上时），我们有
\begin{equation*}
\mathrm{i}G(t, 0) = \int \frac{\mathrm{d}^{d}\pmb{k}}{(2\pi)^{d}}\left(+\Theta(t)\Theta(\xi_{\pmb{k}})\mathrm{e}^{-\mathrm{i}t\xi_{\pmb{k}}} - \Theta(-t)\Theta(-\xi_{\pmb{k}})\mathrm{e}^{-\mathrm{i}t\xi_{\pmb{k}}}\right)
\end{equation*}

%FIG{4.5}: {可用于隧穿的态的数目正比于 $V$。}

上式可以通过引入态密度（density of states，单位体积单位能量内态的数目）
\begin{equation*}
N(\epsilon) = \int \frac{\mathrm{d}^{d}\pmb{k}}{(2\pi)^{d}}\,\delta(\epsilon_{\pmb{k}} - \epsilon)
\end{equation*}
来计算，其中
\begin{align*}
N(\epsilon) &= \sqrt{\frac{m}{2\pi^{2}\epsilon}} && \text{当 } d = 1\\
N(\epsilon) &= \frac{m}{2\pi} && \text{当 } d = 2\\
N(\epsilon) &= \frac{m\sqrt{2m\epsilon}}{2\pi^{2}} && \text{当 } d = 3 \tag{4.2.12}
\end{align*}
在 $T = 0$ 且 $t$ 很大时，只有 $\epsilon = E_{F}$ 附近的 $N(\epsilon)$ 才是重要的。于是我们可以假定 $N(\epsilon) = N(E_{F})$ 为常数。代入 $\int \mathrm{d}^{d}\pmb{k}/(2\pi)^{d} = \int \mathrm{d}\xi\, N(\xi + E_{F})$ 后，对于大的 $t$，我们求得
\begin{equation*}
G(t, 0) = -\frac{N(E_{F})}{t - \mathrm{i}0^{+}\operatorname{sgn}(t)}
\end{equation*}
长时间上的代数关联是被占据态的态密度中不连续性的后果。如果态密度按 $N(\epsilon) \propto |\epsilon - E_{F}|^{g}$ 的方式趋于零（见图 4.4），则长时间关联将衰减得更快（假定 $g > 0$）：
\begin{equation*}
G(t, 0) \propto -\mathrm{i}\operatorname{sgn}(t)\,\mathrm{e}^{-\mathrm{i}\frac{\pi}{2}(1+g)}\frac{1}{|t|^{1+g}}
\end{equation*}

费米子格林函数只有在实验中能够测量才有意义。为了考察如何在实验中测量费米子格林函数，考虑两个态密度为 $N_{R,L}(\xi) \propto |\xi|^{g_{R,L}}$ 的金属之间的隧穿（tunneling）。（注意，对普通金属有 $g_{R,L} = 0$。）假定隧穿哈密顿量为
\begin{equation*}
H_{T} = \Gamma(\mathrm{e}^{\mathrm{i}Vt}c_{R}^{\dagger}c_{L} + \text{h.c.})
\end{equation*}
则从左到右的隧穿电流为（见式 (3.7.15)）
\begin{align*}
\langle j_{T}\rangle(t) &= -\Gamma^{2}\int^{t}\mathrm{d}t'\left(\mathrm{e}^{\mathrm{i}V(t-t')}\left\langle [I(t), I^{\dagger}(t')]\right\rangle + \text{h.c.}\right)\\
&= (-\mathrm{i}\Gamma^{2}D^{\beta}(V) + \text{h.c.}) = 2\Gamma^{2}\mathrm{Im}D^{\beta}(V)
\end{align*}
其中 $I(t) = c_{R}^{\dagger}(t,0)c_{L}(t,0)$ 是隧穿算符。隧穿算符 $I(t) = c_{R}^{\dagger}(t,0)c_{L}(t,0)$ 的关联为（设 $t > 0$ 且 $T = 0$）
\begin{align*}
\left\langle I(t)I^{\dagger}(0)\right\rangle &= \left\langle c_{L}(t,0)c_{L}^{\dagger}(0,0)\right\rangle\left\langle c_{R}^{\dagger}(t,0)c_{R}(0,0)\right\rangle\\
&= -G_{L}(t,0)(-)G_{R}(-t,0)\\
&\propto \mathrm{e}^{-\mathrm{i}\frac{\pi}{2}(2+g_{R}+g_{L})}\frac{1}{t^{2+g_{R}+g_{L}}}
\end{align*}
我们看到，隧穿实验测量的是结两侧等空间格林函数的乘积。上述关联与两个 $(1+1)$ 维相互作用玻色系统之间隧穿算符的关联具有相同的形式。重复 3.7.4 节中的计算，我们求得隧穿电流为
\begin{equation*}
I \propto V^{1+g_{R}+g_{L}} = V \times V^{g_{R}} \times V^{g_{L}}
\end{equation*}
当 $g_{L} = g_{R} = 0$ 时，很容易看出可用于隧穿的态的数目正比于 $V$（图 4.5），因而 $I \propto V$。当 $g_{L}$ 与 $g_{R}$ 不为零时，可用的态将受到相应的压制，因而隧穿电流被因子 $V^{g_{R}} \times V^{g_{L}}$ 压低。

我们想强调，隧穿直接测量的是费米子格林函数。虽然上面我们只讨论了自由费米子系统，这些结果同样适用于相互作用的费米子。设想相互作用改变了费米子格林函数中的长时间衰减指数 $g$；那么这一改变就可以通过隧穿实验测出。

\subsection{费米子谱函数}

\begin{keypoints}
\item 费米子格林函数的谱函数（spectral function）。
\item 费米子格林函数（或更确切地说，费米子谱函数的交叠）可以在隧穿实验中测量。
\end{keypoints}

对相互作用的电子，$G_{L,R}(t,0)$ 要更复杂。为了理解两个相互作用费米子系统之间的隧穿，引入费米子格林函数的谱表示（spectral representation）是有用的：\footnote{如果取 $\eta = 1$，下面的讨论同样适用于玻色算符。}
\begin{equation*}
\mathrm{i}G^{\beta}(t,0) = \left\{\begin{array}{ll}
\int \mathrm{d}\omega\, A_{+}^{0}(\omega)\mathrm{e}^{-\mathrm{i}\omega t}, & \text{当 } t > 0\\
\eta \int \mathrm{d}\omega\, A_{-}^{0}(\omega)\mathrm{e}^{-\mathrm{i}\omega t}, & \text{当 } t < 0
\end{array}\right.
\end{equation*}
其中 $A_{+,-}^{0}(\omega)$ 为下列谱函数：
\begin{equation*}
A_{+}^{0}(\nu) = \sum_{m,n}\delta[\nu - (\epsilon_{m} - \epsilon_{n})]\langle \psi_{n}|c(\pmb{x})|\psi_{m}\rangle\langle \psi_{m}|c^{\dagger}(\pmb{x})|\psi_{n}\rangle\frac{\mathrm{e}^{-\epsilon_{n}\beta}}{Z}
\end{equation*}
\begin{equation*}
A_{-}^{0}(\nu) = \sum_{m,n}\delta[\nu + (\epsilon_{m} - \epsilon_{n})]\langle \psi_{n}|c^{\dagger}(\pmb{x})|\psi_{m}\rangle\langle \psi_{m}|c(\pmb{x})|\psi_{n}\rangle\frac{\mathrm{e}^{-\epsilon_{n}\beta}}{Z}
\end{equation*}
并且，对费米子算符 $\eta = -1$（如果我们想为玻色子格林函数引入谱表示，则取 $\eta = 1$）。我们也可以在动量空间中引入谱函数：
\begin{equation*}
\mathrm{i}G^{\beta}(t,\pmb{k}) = \left\{\begin{array}{ll}
\int \mathrm{d}\omega\, A_{+}(\omega,\pmb{k})\mathrm{e}^{-\mathrm{i}\omega t} & \text{当 } t > 0\\
\eta \int \mathrm{d}\omega\, A_{-}(\omega,\pmb{k})\mathrm{e}^{-\mathrm{i}\omega t} & \text{当 } t < 0
\end{array}\right. \tag{4.2.13}
\end{equation*}
其中
\begin{equation*}
A_{+}(\nu,\pmb{k}) = \sum_{m,n}\delta[\nu - (\epsilon_{m} - \epsilon_{n})]\langle \psi_{n}|c_{\pmb{k}}|\psi_{m}\rangle\langle \psi_{m}|c_{\pmb{k}}^{\dagger}|\psi_{n}\rangle\frac{\mathrm{e}^{-\epsilon_{n}\beta}}{Z}
\end{equation*}
\begin{equation*}
A_{-}(\nu,\pmb{k}) = \sum_{m,n}\delta[\nu + (\epsilon_{m} - \epsilon_{n})]\langle \psi_{n}|c_{\pmb{k}}^{\dagger}|\psi_{m}\rangle\langle \psi_{m}|c_{\pmb{k}}|\psi_{n}\rangle\frac{\mathrm{e}^{-\epsilon_{n}\beta}}{Z}
\end{equation*}
我们看到
\begin{equation*}
A_{\pm}^{0}(\omega) = \mathcal{V}^{-1}\sum_{\pmb{k}}A_{\pm}(\omega,\pmb{k})
\end{equation*}
由定义我们还看到，$A_{+,-}^{0}$ 与 $A_{+,-}$ 是实的且为正。在 $\omega$--$\pmb{k}$ 空间中，我们求得
\begin{align*}
G_{+}^{\beta}(\omega,\pmb{k}) &\equiv \int \mathrm{d}t\,\Theta(t)G^{\beta}(t,\pmb{k})\mathrm{e}^{\mathrm{i}\omega t} = \int \mathrm{d}\nu\,\frac{A_{+}(\nu,\pmb{k})}{\omega - \nu + \mathrm{i}0^{+}}\\
G_{-}^{\beta}(\omega,\pmb{k}) &\equiv \int \mathrm{d}t\,\Theta(-t)G^{\beta}(t,\pmb{k})\mathrm{e}^{\mathrm{i}\omega t} = -\eta \int \mathrm{d}\nu\,\frac{A_{-}(\nu,\pmb{k})}{\omega - \nu - \mathrm{i}0^{+}} \tag{4.2.14}
\end{align*}
动量空间中的谱函数 $A_{+,-}(\nu,\pmb{k})$ 完全决定了格林函数。

%FIG{4.6}: {谱函数及其交叠。}

响应函数 $\langle [I(t), I^{\dagger}(0)]\rangle$ 可以用 R、L 两种金属中的谱函数表出。对 $t > 0$，我们求得
\begin{align*}
&\left\langle [I(t), I^{\dagger}(0)]\right\rangle\\
={}& \left\langle c_{L}(t,0)c_{L}^{\dagger}(0,0)\right\rangle\left\langle c_{R}^{\dagger}(t,0)c_{R}(0,0)\right\rangle - \left\langle c_{L}^{\dagger}(0,0)c_{L}(t,0)\right\rangle\left\langle c_{R}(0,0)c_{R}^{\dagger}(t,0)\right\rangle\\
={}& (\mathrm{i}G_{L}(t))(-\mathrm{i}G_{R}(-t)) - (-\mathrm{i}G_{L}(-t))^{*}(\mathrm{i}G_{R}(t))^{*} \tag{4.2.15}\\
={}& -\int \mathrm{d}\omega_{L}\mathrm{d}\omega_{R}\left(A_{L+}^{0}(\omega_{L})A_{R-}^{0}(\omega_{R}) - A_{L-}^{0}(\omega_{L})A_{R+}^{0}(\omega_{R})\right)\mathrm{e}^{\mathrm{i}(\omega_{R}-\omega_{L})t}
\end{align*}
于是，响应函数 $\mathrm{i}D^{\beta}(t) = \Theta(t)[I(t), I^{\dagger}(0)]$ 在 $\omega$ 空间中由下式给出：
\begin{align*}
D^{\beta}(\omega) = {}& \int \mathrm{d}t\, D^{\beta}(t)\mathrm{e}^{\mathrm{i}\omega t} \tag{4.2.16}\\
={}& \int \mathrm{d}\omega_{L}\mathrm{d}\omega_{R}\,\frac{A_{L+}^{0}(\omega_{L})A_{R-}^{0}(\omega_{R}) - A_{L-}^{0}(\omega_{L})A_{R+}^{0}(\omega_{R})}{\omega - \omega_{L} + \omega_{R} + \mathrm{i}0^{+}}
\end{align*}
从左到右的隧穿电流由虚部确定：
\begin{align*}
\langle j_{T}\rangle(t) ={}& 2\Gamma^{2}\mathrm{Im}D^{\beta}(V)\\
={}& 2\pi\Gamma^{2}\int \mathrm{d}\nu\left(A_{L-}^{0}(\nu)A_{R+}^{0}(\nu - V) - A_{L+}^{0}(\nu)A_{R-}^{0}(\nu - V)\right) \tag{4.2.17}
\end{align*}
在零温下，$A_{+}(\omega)$ 仅当 $\omega > 0$ 时非零，而 $A_{-}(\omega)$ 仅当 $\omega < 0$ 时非零。因此，两式中只有一项有贡献，取决于 $V$ 的符号（见图 4.6）。该贡献是隧穿结（tunneling junction）两侧谱函数的一个交叠积分（overlap integral）。

对自由费米子，由虚时格林函数 (4.2.9) 出发，我们求得（见式 (2.2.37)）
\begin{equation*}
A_{-}(\omega,\pmb{k}) = n_{F}(\xi_{\pmb{k}})\delta(\omega - \xi_{\pmb{k}}), \qquad A_{+}(\omega,\pmb{k}) = (1 - n_{F}(\xi_{\pmb{k}}))\delta(\omega - \xi_{\pmb{k}})
\end{equation*}
（把上式代入式 (4.2.14)，便可恢复有限温度下由式 (4.2.5) 给出的自由电子格林函数。）对 $\pmb{k}$ 积分后，我们得到
\begin{equation*}
A_{-}^{0}(\omega) = n_{F}(\omega)N(\omega + E_{F}), \qquad A_{+}^{0}(\omega) = (1 - n_{F}(\omega))N(\omega + E_{F}) \tag{4.2.18}
\end{equation*}
隧穿电流为
\begin{align*}
\langle j_{T}\rangle(t) = {}& 2\pi\Gamma^{2}\int \mathrm{d}\nu \left[(1 - n_{F}(\nu))N_{L}(\nu + E_{F})n_{F}(\nu - V)N_{R}(\nu - V + E_{F})\right.\\
&\left.- n_{F}(\nu)N_{L}(\nu + E_{F})(1 - n_{F}(\nu - V))N_{R}(\nu - V + E_{F})\right]\\
\approx {}& 2\pi\Gamma^{2}N_{L}(E_{F})N_{R}(E_{F})\int \mathrm{d}\nu\left(n_{F}(\nu - V) - n_{F}(\nu)\right)
\end{align*}
这是一个非常简单的结果，可以直接用二阶微扰论得到。

\subsection*{问题 4.2.5}

计算 $d = 1$ 维时实时空中的编时费米子格林函数在 $(t, x) \to \infty$ 极限下的结果（但 $t/x$ 任意）。在 $t/x \to 0$ 与 $t/x \to \infty$ 两个极限下检验你的结果。

\subsection*{问题 4.2.6}

两个相同的 $d = 1$ 维自由费米子系统在位于 $x = 0$ 和 $x = a$ 的两点处连接。隧穿哈密顿量为
\begin{equation*}
H_{T} = \Gamma\left[\left[c_{L}(0)c_{R}^{\dagger}(0) + c_{L}(a)c_{R}^{\dagger}(a)\right] + \text{h.c.}\right]
\end{equation*}
利用上一题的结果，求作为 $a$ 的函数的隧穿电导。

\subsection*{问题 4.2.7}

证明有限温度下的谱求和规则（spectral sum rule）
\begin{equation*}
\int \mathrm{d}\omega \left(A_{+}(\omega,\pmb{k}) + A_{-}(\omega,\pmb{k})\right) = 1
\end{equation*}
（提示：利用 $\{c_{k}, c_{k}^{\dagger}\} = 1$。）

\subsection*{问题 4.2.8}

\textbf{平行二维电子系统之间的隧穿}（Eisenstein et al., 1991）：考虑二维的两个相同费米子系统，由一个竖直面积隧穿结（vertical area tunneling junction）连接，该结保持二维动量守恒（见图 4.7(a)）。隧穿哈密顿量具有如下形式：

%FIG{4.7}: {(a) 竖直面积隧穿结；(b) 两张发生错移的费米面。}

\begin{equation*}
H_{T} = A \sum_{\pmb{k}}\left[c_{L\pmb{k}}c_{R\pmb{k}}^{\dagger} + \text{h.c.}\right]
\end{equation*}
在没有平行于二维层的磁场时，两层中的费米子具有相同的色散关系 $\epsilon_{\pmb{k}}$。当存在平行于二维层的磁场 $B$ 时，两层中费米子的动量彼此相对错移（见图 4.7(b)），两层中的色散变为
\begin{equation*}
\epsilon_{R\pmb{k}} = \epsilon_{\pmb{k} - \frac{1}{2}\pmb{Q}}, \qquad \epsilon_{L\pmb{k}} = \epsilon_{\pmb{k} + \frac{1}{2}\pmb{Q}}
\end{equation*}
其中 $Q \propto B$ 且 $\pmb{Q} \cdot \pmb{B} = 0$。

\begin{enumerate}
\item 证明：对自由费米子，当 $B = 0$ 且电压有限（$V \neq 0$）时，有限温度下的隧穿电流为零。
\item 假定由于相互作用，费米子的谱函数具有如下形式：
\begin{equation*}
A_{+}(\omega,\pmb{k}) = (1 - n_{F}(\omega))\frac{\pi^{-1}\Gamma}{(\omega - \xi_{\pmb{k}})^{2} + \Gamma^{2}}, \qquad A_{-}(\omega,\pmb{k}) = n_{F}(\omega)\frac{\pi^{-1}\Gamma}{(\omega - \xi_{\pmb{k}})^{2} + \Gamma^{2}}
\end{equation*}
其中 $\Gamma \ll E_{F}$ 是一个衰变率（注意上述 $A_{\pm}$ 满足式 (2.2.28)）。求温度小而非零（$T \ll E_{F}$）时隧穿电流的表达式。
\item 当 $B = 0$ 时，在 $T \to 0$ 与 $T \gg \Gamma$ 极限下，隧穿电导对温度的依赖关系是什么？估算 $T = 0$ 时单位面积的隧穿电导。
\item 假定零温下 $B = 0$ 时的隧穿电导为 $\sigma_{0}$，试以 $\sigma_{0}$ 估算有限 $Q$ 下的隧穿电导 $\sigma_{Q}$。当 $Q$ 靠近 0 以及靠近 $2k_{F}$ 时，$\sigma_{Q}$ 的行为如何？
\end{enumerate}

\subsection{等时格林函数与费米面的形状}

\begin{keypoints}
\item 长距离上的代数衰减与动量空间中费米子占据数的锐利特征（不连续性）相关。
\end{keypoints}

我们转向等时（equal-time）格林函数。等时关联完全由基态波函数决定。从定义也容易看出，$G^{\beta}(0^{+},\pmb{x})$ 与 $G^{\beta}(-0^{+},\pmb{x})$ 相差一个反对易子（若算符是玻色型的，则为对易子）的平均：
\begin{equation*}
G^{\beta}(0^{+},\pmb{x}) - G^{\beta}(-0^{+},\pmb{x}) = \left\langle \{c(\pmb{x}), c^{\dagger}(0)\}\right\rangle = \delta(\pmb{x})
\end{equation*}
这里我们想选定一个方向 $\hat{\pmb{x}}$，考察 $G(\pm 0^{+}, x\hat{\pmb{x}})$ 在大 $x$ 下的行为。

在 $T = 0$ 时，$G(-0^{+},\pmb{x})$ 可以写为
\begin{equation*}
\mathrm{i}G(-0^{+},\pmb{x}) = -\int \frac{\mathrm{d}^{d}\pmb{k}}{(2\pi)^{d}}\, n_{F}(\xi_{\pmb{k}})\mathrm{e}^{\mathrm{i}\pmb{k}\cdot\pmb{x}} = -\int_{-\infty}^{+\infty}\mathrm{d}k\, \tilde{N}(k,\hat{\pmb{x}})\mathrm{e}^{\mathrm{i}k|\pmb{x}|}
\end{equation*}
其中
\begin{equation*}
\tilde{N}(k,\hat{\pmb{x}}) = \int \frac{\mathrm{d}^{d}\pmb{k}}{(2\pi)^{d}}\,\Theta(-\xi_{\pmb{k}})\delta(k - \pmb{k}\cdot\hat{\pmb{x}})
\end{equation*}
它可以视为动量空间中被占据态的密度。

%FIG{4.8}: {(a) $\pmb{k}_{F}(\hat{\pmb{x}})$ 的定义，其中的直线与 $x$ 垂直。(b) 带有一块平坦区域的费米面。}

从图 4.8(a) 我们看到，一般地，$\tilde{N}(k,\hat{\pmb{x}})$ 在 $k = \pmb{k}_{F}(\hat{\pmb{x}})\cdot\hat{\pmb{x}}$ 与 $k = \pmb{k}_{F}(-\hat{\pmb{x}})\cdot\hat{\pmb{x}}$ 处含有两个奇点：
\begin{align*}
\tilde{N}(k,\hat{\pmb{x}}) ={}& c_{+}\Theta(\pmb{k}_{F}(\hat{\pmb{x}})\cdot\hat{\pmb{x}} - k)\,|\pmb{k}_{F}(\hat{\pmb{x}})\cdot\hat{\pmb{x}} - k|^{(d-1)/2}\\
&+ c_{-}\Theta(-\pmb{k}_{F}(-\hat{\pmb{x}})\cdot\hat{\pmb{x}} + k)\,\left|-\pmb{k}_{F}(-\hat{\pmb{x}})\cdot\hat{\pmb{x}} + k\right|^{(d-1)/2}
\end{align*}
（例如在一维情形，$\tilde{N}(k,\hat{\pmb{x}})$ 有两个不连续的台阶。）这两个奇点导致如下代数长程的等时关联：
\begin{equation*}
\mathrm{i}G(-0^{+},\pmb{x})\big|_{\pmb{x}\to\infty} \sim \left(c_{+}\mathrm{e}^{\mathrm{i}\pmb{k}_{F}(\hat{\pmb{x}})\cdot\pmb{x}}\mathrm{e}^{-\mathrm{i}\frac{\pi(d+1)}{4}} + c_{-}\mathrm{e}^{\mathrm{i}\pmb{k}_{F}(-\hat{\pmb{x}})\cdot\pmb{x}}\mathrm{e}^{+\mathrm{i}\frac{\pi(d+1)}{4}}\right)\frac{1}{|\pmb{x}|^{(d+1)/2}}
\end{equation*}
对球形的费米面，我们有
\begin{equation*}
\mathrm{i}G(-0^{+},\pmb{x})\big|_{\pmb{x}\to\infty} \sim \cos\left(k_{F}|\pmb{x}| - \frac{\pi(d+1)}{4}\right)\frac{1}{|\pmb{x}|^{(d+1)/2}}
\end{equation*}
如果费米面包含一块平坦区域（见图 4.8(b)），那么 $\tilde{N}(k,\hat{\pmb{x}})$ 在相应方向上会有一个不连续台阶。在该方向上，$G(-0^{+},\pmb{x})$ 将具有较慢的代数衰减，其衰减幂次与一维系统的相同：
\begin{equation*}
\mathrm{i}G(-0^{+},\pmb{x})\big|_{\pmb{x}\to\infty} \sim -\mathrm{i}\,\mathrm{e}^{\mathrm{i}\pmb{k}_{F}(\hat{\pmb{x}})\cdot x}\frac{1}{|\pmb{x}|}
\end{equation*}

\section{两体关联函数与线性响应}

\begin{keypoints}
\item 我们可以由两体关联函数计算线性响应。这使我们能够计算许多可测量的量。
\item 两体关联函数揭示自由费米子基态的许多潜在不稳定性。
\end{keypoints}

两体关联函数（two-body correlation functions）是最重要的一类关联。实验中测量的大多数物理量都直接与两体关联相关。这些物理量包括电导与热导、中子与光散射截面、弹性常数、介电常数、磁化率等等（见 3.7.2 节）。下面我们将讨论无自旋自由费米子系统的两体关联。

\subsection{密度--密度关联函数}

我们首先考虑两体关联中的一种——密度关联函数
\begin{equation*}
\mathrm{i}P^{00}(t,\pmb{x}) = \langle T(\rho(t,\pmb{x})\rho(0))\rangle
\end{equation*}
其中
\begin{equation*}
\rho(t,\pmb{x}) = c^{\dagger}(t,\pmb{x})c(t,\pmb{x}).
\end{equation*}
密度关联函数与压缩率以及中子/光散射相关。为了计算 $P^{00}(t,\pmb{x})$，我们可以对自由费米子算符应用 Wick 定理，如下：

%FIG{4.9}: {具有嵌套波矢 $Q$ 的嵌套费米面。}

\begin{align*}
\mathrm{i}P^{00}(t,\pmb{x}) ={}& \left\langle T[c^{\dagger}(t + 0^{+},\pmb{x})c(t,\pmb{x})c^{\dagger}(0^{+})c(0)]\right\rangle\\
={}& \rho_{0}^{2} - \mathrm{i}G(-t,-\pmb{x})\,iG(t,\pmb{x})
\end{align*}
利用前面已算出的费米子格林函数，我们发现，在 $t = 0$ 极限下，对于球形费米面，有
\begin{equation*}
\mathrm{i}P^{00}(0,\pmb{x}) = \rho_{0}^{2} + C\left(1 - \sin\left(2k_{F}|\pmb{x}| - \frac{\pi(d+1)}{2}\right)\right)\frac{1}{|\pmb{x}|^{(d+1)}}
\end{equation*}
对图 4.9 中的嵌套费米面（nested Fermi surface），我们有
\begin{equation*}
\mathrm{i}P^{00}(0,\pmb{x}) = \rho_{0}^{2} + C[1 + \sin(Q|\pmb{x}|)]\frac{1}{|\pmb{x}|^{2}}
\end{equation*}
其中 $x \| Q$。

我们知道，在晶体中，密度关联会一直振荡而没有任何衰减，即 $\mathrm{i}P^{00}(0,\pmb{x}) = \rho_{0}^{2} + C[1 + \sin(Q|\pmb{x}|)]$。振荡部分代表了有序晶体中的长程序。我们看到，自由的费米子基态没有长程晶体序，但它具有“代数长程”晶体序，对嵌套费米面尤其如此（见图 4.10）。在某种意义上，嵌套的自由费米子系统正处于成为晶体的边缘。我们将在本章后面看到，只要开启一个小的相互作用，“代数长程”晶体序就可以被提升为长程晶体序。
为了得到自由费米子的线性响应，我们需要计算密度响应函数 $\Pi^{00}$。在 $t$--$\pmb{k}$ 空间中，密度响应函数可以
